\chapter*{Abstract}
\addcontentsline{toc}{chapter}{Abstract}

This work presents an optimal route management system for parcel delivery in the city of Cusco, implemented as a Python desktop application with a Tkinter graphical interface. The road network is modelled as a directed graph of 20 points of interest and 64 connections weighted by real street lengths from OpenStreetMap; shortest distances between places are precomputed with a hand-written Dijkstra algorithm. On top of this model, the four required paradigms were implemented without external optimization libraries: (1) \textit{Greedy}, with the nearest-neighbour heuristic and the assignment of urgent orders to the closest courier; (2) \textit{Divide and Conquer}, which splits the city at the median alternating latitude and longitude, both into zones of $K$ orders and into one zone per courier; (3) \textit{Dynamic Programming}, with the 0/1 Knapsack to load the vehicle and Held-Karp for the exact tour; and (4) \textit{Backtracking}, which enumerates every simple path between two places while avoiding closed streets, with optional bound-based pruning. Measurements confirmed the Big-O analysis: Held-Karp took 22.2 ms for $n=12$ points against 0.017 ms for Greedy, while the greedy tour was up to 24.5\% longer than the optimum. In the case stated by the assignment, from San Blas to Wanchaq, Backtracking found all 27 simple paths and the shortest one (1,778.6 m); pruning kept that same path while visiting 20 states instead of 146. Splitting the 30 orders into three zones, one per courier, reduced the total distance from 43.4 km (a single greedy courier) to 30.2 km. The system passes 39 unit tests and an automated interface test that exercises every feature.

\textit{Keywords:} optimal routes, greedy algorithms, divide and conquer, dynamic programming, backtracking, Held-Karp, 0/1 knapsack, Dijkstra, Cusco
